% SHARD-ID: AQM-26-PRESHEAVES-COSHEAVES-OBSERVABLE-NETS
% SHARD-TITLE: Presheaves, Cosheaves, and Observable Nets
% SHARD-SUMMARY: Defines presheaves, sheaves, precosheaves, and cosheaves before applying them to the observable net.
% SHARD-SUMMARY: Proves that the open-local observable assignment is a precosheaf but not an ordinary cosheaf of unital \(*\)-algebras.
% SHARD-SUMMARY: Shows that the Weyl label assignment is the cosheaf-like layer whose quantization gives additive local nets.
% SHARD-KEYWORDS: presheaf, sheaf, precosheaf, copresheaf, cosheaf, local net, observable algebra, Weyl labels

\section{Presheaves, Cosheaves, and Observable Nets}
\label{sec:presheaves-cosheaves-observable-nets}

\subsection{Status and Source Anchors}

Stacks Project
\path{references/algebraic_geometry/stacks_project_sheaves.tex} defines
presheaves in lines 73--123, sheaves of sets in lines 500--525, and
category-valued sheaves by an equalizer diagram in lines 719--736.  It also
records that \(U\mapsto\bigoplus_{x\in U}M_x\) is generally not a sheaf in
lines 622--637.

Curry's cosheaf source
\path{references/cosheaves/curry_1303_3255_source/ch_abstract_sheaves.tex}
defines pre-sheaves and pre-cosheaves in lines 89--92, sheaves and cosheaves
by cover limits and colimits in lines 102--118, and rewrites the cosheaf
condition as a coequalizer in lines 184--218.  The vector-space exact sequence
form is in lines 243--259.  Curry's category source
\path{references/cosheaves/curry_1303_3255_source/ch_categories.tex} defines
colimits in lines 337--358 and records coproducts and unions as colimits in
lines 366--384.

\subsection{Lamport Dependency Skeleton}

\[
\begin{array}{rcl}
D1&:&\mathsf O(X)\text{ is the category of open subsets ordered by inclusion.}\\
D2&:&\text{A presheaf is a functor }\mathsf O(X)^{op}\to\mathcal C.\\
D3&:&\text{A sheaf is a presheaf satisfying cover equalizer gluing.}\\
D4&:&\text{A precosheaf is a functor }\mathsf O(X)\to\mathcal C.\\
D5&:&\text{A cosheaf is a precosheaf satisfying cover coequalizer gluing.}\\
D6&:&U\mapsto E(U)=\bigoplus_{x\in U}E_x\text{ is the label precosheaf.}\\
D7&:&U\mapsto\mathcal A(U)\text{ is the open-local observable assignment.}\\
L1&:&E\text{ satisfies the finite-cover cosheaf axiom in abelian groups.}\\
L2&:&\mathcal A\text{ is a precosheaf of unital }*\text{-algebras.}\\
L3&:&\mathcal A\text{ satisfies net additivity inside the ambient algebra.}\\
T1&:&\mathcal A\text{ is not an ordinary cosheaf of unital }*\text{-algebras.}\\
T2&:&\text{Quantization turns label cosheaf gluing into tensor-local net gluing.}
\end{array}
\]

\subsection{Presheaves and Sheaves D1--D3}

Let \(X\) be a topological space.  The category \(\mathsf O(X)\) has open
sets \(U\subset X\) as objects and a unique morphism \(U\to V\) when
\(U\subset V\).

Let \(\mathcal C\) be a category.  A \(\mathcal C\)-valued presheaf on \(X\)
is a contravariant functor
\[
  F:\mathsf O(X)^{op}\longrightarrow\mathcal C.
\]
Equivalently, for \(V\subset U\) it has restriction morphisms
\[
  \rho^U_V:F(U)\longrightarrow F(V)
\]
with \(\rho^U_U=\operatorname{id}\) and
\(\rho^U_W=\rho^V_W\rho^U_V\) for \(W\subset V\subset U\).

Assume \(\mathcal C\) has the products and equalizers needed for the displayed
cover diagram.  A presheaf \(F\) is a sheaf if for every open cover
\(U=\bigcup_iU_i\) the diagram
\[
  F(U)\longrightarrow \prod_i F(U_i)
  \rightrightarrows \prod_{i,j}F(U_i\cap U_j)
\]
is an equalizer.  In sets this says that compatible local sections glue to a
unique section over \(U\).

\subsection{Precosheaves and Cosheaves D4--D5}

We use ``precosheaf'' and ``copresheaf'' synonymously.  A
\(\mathcal C\)-valued precosheaf on \(X\) is a covariant functor
\[
  G:\mathsf O(X)\longrightarrow\mathcal C.
\]
For \(V\subset U\) it has extension morphisms
\[
  r_{U,V}:G(V)\longrightarrow G(U)
\]
with \(r_{U,U}=\operatorname{id}\) and
\(r_{U,W}=r_{U,V}r_{V,W}\) for \(W\subset V\subset U\).

Assume \(\mathcal C\) has the coproducts and coequalizers needed for the
displayed cover diagram.  A precosheaf \(G\) is a cosheaf if for every open
cover \(U=\bigcup_iU_i\) the diagram
\[
  \coprod_{i,j}G(U_i\cap U_j)
  \rightrightarrows
  \coprod_iG(U_i)
  \longrightarrow
  G(U)
\]
is a coequalizer.  In an additive category such as abelian groups or vector
spaces, the finite-cover form is the exact sequence
\[
  \bigoplus_{i,j}G(U_i\cap U_j)
  \longrightarrow
  \bigoplus_iG(U_i)
  \longrightarrow
  G(U)\longrightarrow0.
\]
The arrows identify data that have been counted twice on overlaps.

\subsection{The Label Cosheaf D6--L1}

At the finite residue-qudit level, each point \(x\in X\) has a local Weyl
label group \(E_x=\kappa(x)^2\), regarded as an abelian group.  Define
\[
  E(U)=\bigoplus_{x\in U}E_x.
\]
For \(V\subset U\), the extension map \(E(V)\to E(U)\) is extension by zero
outside \(V\).

\paragraph{Lemma \(L1\).}
For finite covers, \(E\) satisfies the cosheaf axiom in abelian groups.

\emph{Proof.}
Let \(U=\bigcup_iU_i\) be a finite open cover.  The map
\[
  \bigoplus_iE(U_i)\longrightarrow E(U)
\]
adds the extension-by-zero images.  It is surjective because every summand
\(E_x\) with \(x\in U\) lies in at least one \(U_i\).

The only relations needed are overlap relations.  If \(x\in U_i\cap U_j\),
then the two copies of \(E_x\) inside \(E(U_i)\oplus E(U_j)\) have the same
image in \(E(U)\), and their difference is exactly in the image of
\(E(U_i\cap U_j)\).  Conversely, if a finite sum maps to zero in \(E(U)\),
then for each point \(x\) the sum of its finitely many copies over the index
set \(\{i:x\in U_i\}\) is zero.  This index set is connected through overlaps,
because every pair \(i,j\) in it has \(x\in U_i\cap U_j\).  Hence the
pointwise zero relation is generated by the pairwise overlap relations.  Since
the direct sum over points separates these pointwise calculations, the
displayed sequence is exact.

\subsection{The Observable Precosheaf D7--L3}

Let \(\mathcal A(U)\) be the open-local observable algebra from AQM-24.  For
\(V\subset U\), define
\[
  \iota_{V,U}:\mathcal A(V)\hookrightarrow\mathcal A(U),
  \qquad a\mapsto a\otimes\mathbf 1_{U\setminus V}.
\]

\paragraph{Lemma \(L2\).}
\(U\mapsto\mathcal A(U)\) is a precosheaf of unital \(*\)-algebras.

\emph{Proof.}
The identity and composition laws are exactly Lemma \(L1\) of AQM-24:
\[
  \iota_{U,U}=\operatorname{id},\qquad
  \iota_{V,W}\iota_{U,V}=\iota_{U,W}
\]
for \(U\subset V\subset W\).  Thus \(\mathcal A\) is a covariant functor from
\(\mathsf O(X)\) to unital \(*\)-algebras.

\paragraph{Lemma \(L3\).}
For finite unions,
\[
  \mathcal A(U\cup V)
  =
  \mathcal A(U)\vee\mathcal A(V),
\]
where \(\vee\) denotes the unital \(*\)-subalgebra generated by the two images
inside \(\mathcal A(U\cup V)\).  If \(U\cap V=\emptyset\), the two images
commute and the generated algebra is the tensor product.

\emph{Proof.}
Every tensor factor of \(\mathcal A(U\cup V)\) belongs to \(U\) or to \(V\).
Hence the images of \(\mathcal A(U)\) and \(\mathcal A(V)\) generate all local
matrix factors in the union.  If \(U\cap V=\emptyset\), the two images act on
disjoint tensor factors, so elementary tensors commute factor by factor.  The
generated algebra is then the tensor product of the two local algebras.

\subsection{Why This Is Not an Ordinary Cosheaf T1}

\paragraph{Theorem \(T1\).}
The observable precosheaf \(\mathcal A\) is not, in general, a cosheaf valued
in the category of unital \(*\)-algebras with ordinary categorical colimits.

\emph{Proof.}
It is enough to use \(X=\operatorname{Spec}(\mathbb Z/6\mathbb Z)\).  Let
\[
  U=\{(2)\},\qquad V=\{(3)\}.
\]
Then \(U\cap V=\emptyset\), \(U\cup V=X\), and AQM-24 gives
\[
  \mathcal A(U)=M_2,\qquad
  \mathcal A(V)=M_3,\qquad
  \mathcal A(X)=M_2\otimes M_3.
\]
Since \(\mathcal A(\emptyset)=\mathbb C\), the cosheaf coequalizer for this
cover is the coproduct of unital \(*\)-algebras amalgamated over the common
unit:
\[
  M_2 *_\mathbb C M_3.
\]
This coproduct is not \(M_2\otimes M_3\).  Indeed, if it were the tensor
product with the canonical embeddings, then every pair of unital
\(*\)-homomorphisms \(M_2\to B\) and \(M_3\to B\) would have commuting images,
because the two canonical tensor factors commute.

Take \(B=M_6(\mathbb C)\), written as \(3\times3\) blocks in a \(2\times2\)
matrix.  Let
\[
  f(a)=a\otimes I_3,
  \qquad
  g(b)=
  \begin{bmatrix}
    b&0\\
    0&VbV^*
  \end{bmatrix},
\]
where \(V\in M_3(\mathbb C)\) swaps the first two basis vectors.  These are
unital \(*\)-homomorphisms.  For the matrix units \(e_{12}\in M_2\) and
\(E_{11}\in M_3\),
\[
  f(e_{12})g(E_{11})
  =
  \begin{bmatrix}0&E_{22}\\0&0\end{bmatrix},
  \qquad
  g(E_{11})f(e_{12})
  =
  \begin{bmatrix}0&E_{11}\\0&0\end{bmatrix}.
\]
These are unequal.  Therefore the universal object for arbitrary unital
\(*\)-homomorphisms cannot be the tensor product.  The cosheaf colimit in
unital \(*\)-algebras is too free: it does not impose locality commutation.

\subsection{Quantization and the Correct Dictionary T2}

\paragraph{Theorem \(T2\).}
The correct next-layer dictionary is:
\[
\begin{array}{rcl}
E&:&\text{a cosheaf-like label object in abelian groups or vector spaces},\\
\mathcal A&:&\text{a precosheaf/local net of observable algebras},\\
&&\text{with additivity and locality imposed inside the ambient algebra.}
\end{array}
\]

\emph{Proof.}
Lemma \(L1\) proves that labels glue by the expected additive colimit.  Lemmas
\(L2\) and \(L3\) prove that observable algebras form a covariant local net
and satisfy the physical additivity relation
\[
  \mathcal A(U\cup V)=\mathcal A(U)\vee\mathcal A(V).
\]
Theorem \(T1\) proves that this additivity is not the ordinary cosheaf axiom
in unital \(*\)-algebras: categorical coproducts give free products, whereas
quantized independent labels give commuting tensor factors.  Thus quantization
does not preserve the raw cosheaf colimit from labels to observable algebras;
it turns additive label gluing into tensor-local algebra gluing.

\subsection{Conclusion}

The open-local observable assignment is naturally a precosheaf:
\[
  U\subset V\quad\Rightarrow\quad
  \mathcal A(U)\hookrightarrow\mathcal A(V).
\]
It should not be called an ordinary cosheaf of unital \(*\)-algebras unless the
target category or gluing notion is changed.  The label layer \(E\) is the
cosheaf-like layer.  The observable layer is the AQFT-style net: isotone,
local on disjoint opens, and additive by generated subalgebras.
